\section{Transacciones distribuidas: Two-Phase Commit}

Una compra toca dos servicios: \texttt{ServicioPagos} debita el monto y
\texttt{ServicioInventario} asigna el asiento. Ninguno puede confirmar solo.
El coordinador corre un Two-Phase Commit. El código está en el módulo
\texttt{dos\_fases}. Se corre con \texttt{cargo run -- parte3}.

\subsection{Fase Prepare y fase Commit/Rollback}

La función \texttt{coordinar\_2pc} manda \texttt{Prepare} a los dos
participantes y espera los votos con \texttt{recv\_timeout}. Cada servicio
es un hilo con buzón. Si puede reservar el recurso, contesta \texttt{Yes}
y pasa a \texttt{Preparado}. Si no, contesta \texttt{No} y aborta en local.

\begin{lstlisting}
// FASE 1
for (_, p) in participantes {
    p.send(MsgParticipante::Prepare {
        tx,
        responder_a: tx_votos.clone(),
    });
}
// junta votos hasta el timeout
let todos_yes = participantes
    .iter()
    .all(|(n, _)| votos.get(n) == Some(&Voto::Yes));

// FASE 2
let decision = if todos_yes {
    Decision::Commit
} else {
    Decision::Rollback
};
for (_, p) in participantes {
    match decision {
        Decision::Commit => {
            p.send(MsgParticipante::Commit { tx });
        }
        Decision::Rollback => {
            p.send(MsgParticipante::Rollback { tx });
        }
    }
}
\end{lstlisting}

En Pagos, \texttt{preparar} resta el monto y marca \texttt{retenido} si el
saldo alcanza. En Inventario, marca el asiento retenido si estaba libre.
\texttt{confirmar} deja el débito o la venta definitivos.
\texttt{revertir} devuelve lo retenido. Un voto que no llega a tiempo se
cuenta como \texttt{No}: el coordinador manda \texttt{Rollback} a todos.

\subsection{Partición después del Yes}

El escenario del enunciado es este: los dos servicios ya contestaron
\texttt{Yes} y, antes de que llegue el \texttt{Commit}, la red se parte o
el coordinador se cae. En el código eso es la falla
\texttt{CaeTrasPrepare}. El coordinador junta los votos, no decide y no
manda nada. Cada participante preparado espera la decisión con
\texttt{recv\_timeout}. Si vence, marca \texttt{bloqueado} y sigue
reteniendo el recurso. No puede confirmar solo: el otro pudo haber votado
\texttt{No}. No puede abortar solo: el otro pudo haber recibido un
\texttt{Commit}. Queda parado con el saldo o el asiento tomados.

Eso es el bloqueo de 2PC. Safety sigue en pie: no hay un servicio en
\texttt{Confirmado} y otro en \texttt{Abortado} por la misma compra. Lo
que se pierde es liveness: la transacción no termina mientras dure la
partición. En términos de CAP, 2PC elige consistencia frente a la
partición y deja de responder sobre esos recursos. Es un diseño CP: la
red no es confiable, así que la elección queda entre C y A, y acá se
prioriza que los dos nodos vean el mismo desenlace.

\subsection{Pruebas}

El test del caso feliz pide \texttt{Commit}, saldo 950, asiento vendido y
nadie bloqueado.

Los dos tests de un \texttt{No} (inventario ocupado o saldo corto) fuerzan
\texttt{Rollback}. El que había reservado revierte.

Si Pagos no contesta, el coordinador cuenta el timeout como \texttt{No} y
manda \texttt{Rollback}. Inventario libera el asiento.

El test de la partición deja la decisión en \texttt{None}. Los dos quedan
en \texttt{Preparado}, con \texttt{bloqueado} en verdadero y el recurso
retenido.

Corrimos \texttt{cargo test --lib dos\_fases}. Las cinco pasan.
